% Generated draft: requires mathematical and source review.
\documentclass[11pt]{article}
\usepackage[margin=25mm]{geometry}
\usepackage{fontspec}
\setmainfont{Latin Modern Roman}
\usepackage{amsmath,amssymb,mathtools,unicode-math}
\setmathfont{Latin Modern Math}
\usepackage{enumitem,needspace,xurl,hyperref}
\setlength{\parindent}{0pt}
\setlength{\parskip}{5pt}
\setlength{\emergencystretch}{3em}
\newcommand{\N}{\mathbb{N}}
\newcommand{\Z}{\mathbb{Z}}
\newcommand{\Q}{\mathbb{Q}}
\newcommand{\R}{\mathbb{R}}
\newcommand{\C}{\mathbb{C}}
\newcommand{\F}{\mathbb{F}}
\newcommand{\A}{\mathbb{A}}
\newcommand{\PP}{\mathbb{P}}
\newcommand{\E}{\mathbb{E}}
\newcommand{\eps}{\varepsilon}
\newcommand{\dd}{\,\mathrm{d}}
\begin{document}
\section*{UnsolvedMath ID: 800972}
\textbf{Source problem code:} catalogue entry
\subsection{Definitions and mathematical statement}
A closed manifold means a compact smooth manifold without boundary. A semi-Riemannian metric is a smooth nondegenerate symmetric bilinear form on each tangent space. Interpret completeness as geodesic completeness: every maximal affinely parametrized geodesic is defined on all of \(\R\). A nontrivial smooth \(1\)-form \(\eta\) means \(\eta\not\equiv 0\). Here a light-like geodesic means a nonconstant affinely parametrized geodesic \(\gamma:\R\to M\) satisfying \(g_{\gamma(t)}(\dot\gamma(t),\dot\gamma(t))=0\) for every \(t\in\R\).

Does there exist a closed manifold \(M\), a geodesically complete semi-Riemannian metric \(g\) on \(M\), and a smooth \(1\)-form \(\eta\not\equiv0\) such that
\[
\forall\,\text{light-like geodesics }\gamma:\R\to M,\quad
\exists\, C_1(\gamma)\in\R\setminus\{0\},\ C_2(\gamma)\in\R,\quad
\forall t\in\R,\qquad
\eta_{\gamma(t)}\bigl(\dot\gamma(t)\bigr)=C_1(\gamma)t+C_2(\gamma)?
\]
The constants may depend on the affinely parametrized geodesic. The displayed signed affine identity is the supplied meaning of growth linearly in both directions.

The source does not specify the dimension or metric signature. If definite metrics are admitted under the term semi-Riemannian, there are no nonconstant light-like geodesics and the universal condition is vacuous. A nonvacuous interpretation therefore requires an indefinite signature, but this requirement is not explicit in the supplied record. The exclusion of constant geodesics and the interpretation of completeness above use the usual geodesic conventions and require review against the source context.
\subsection{Short English statement}
The question asks whether a closed manifold can carry a complete semi-Riemannian metric and a nonzero smooth \(1\)-form whose evaluation on the velocity of every nonconstant affinely parametrized light-like geodesic is affine in time with nonzero slope. The intended signature is unspecified; allowing definite metrics makes the light-like condition vacuous.
\subsection{Sources}
\begin{itemize}
\item Local UnsolvedMath registry: permanent ID \(800972\), exact source problem code \texttt{catalogue entry}, title \emph{Open Problems and Questions about Geodesics --- Question 53}, in the AMR Open Problem Lists collection. The supplied record labels the problem open; this is a supplied status, not an independently verified mathematical assessment.
\item Burns and Matveev, \emph{Open Problems and Questions about Geodesics} (2019), Question 53, as attributed in the supplied record. Supplied source URL: \url{https://arxiv.org/abs/1308.5417}. The record reports access on 2026-10-16 and extraction by \texttt{arxiv\_latex}; these are supplied provenance details, not a claim of independent access or verification.
\end{itemize}
\end{document}
